%!TEX root = forallxyyc.tex

\chapter{关于无障碍性的说明}

我们在制作此 HTML 版本时特别注重无障碍性，特别是对于使用辅助技术（AT）的读者，例如像 \href{https://www.nvaccess.org/download/}{NVDA} 和 \href{https://www.freedomscientific.com/products/software/jaws/}{JAWS} 这样的屏幕阅读器。屏幕阅读器旨在帮助低视力或全盲用户访问和浏览本书这类网页。部分明眼用户有时也偏好使用文本转语音（TTS）软件来朗读网页。相比专为盲人用户设计的屏幕阅读器，TTS 软件功能较少，可能无法将本书全部内容转换为音频。

如果您在使用本书 HTML 版本时配合屏幕阅读器或盲文终端，或者在协助依赖此类工具的学生（例如作为教师或无障碍顾问），并有任何反馈，请\href{mailto:rzach@ucalgary.ca}{与我们联系}！

\paragraph{导航与样式} 我们采用了由 \href{https://vlmantova.github.io/bookml/}{BookML} 提供的 HTML 样式，该样式由 \href{https://eps.leeds.ac.uk/maths/staff/4058/dr-vincenzo-l-mantova}{Vincenzo Mantova} 为利兹大学数学系开发。它提供了可折叠的导航菜单，以及位于每页顶部的按钮，用于选择字体大小、在衬线字体和无衬线字体之间切换，以及选择白底黑字、深色字羊皮纸底或深底浅色等显示选项。

页面正文开头包含一个不可见的链接，可跳过导航侧边栏和样式按钮，直接跳转至正文内容。这些按钮依次是：
\begin{itemize}
  \item “切换侧边栏”，用于打开或关闭导航栏；
  \item “字体设置”，用于设置字体大小、衬线或无衬线字体，以及浅色、羊皮纸色或深色模式；
  \item “下载”，提供下载 PDF 版本的链接；以及
  \item “关于工具栏的信息”，用于弹出一个显示导航键盘快捷键的窗口。
\end{itemize}
这些键盘快捷键是：左右箭头键用于跳转到上一页和下一页，按 `s' 键可切换导航侧边栏。

\paragraph{符号与公式}

本书中的逻辑符号和公式被转换为 \href{https://www.w3.org/Math/}{MathML}，并使用 \href{https://www.mathjax.org/}{MathJax} 显示。MathJax 为公式提供了额外的\href{https://docs.mathjax.org/en/latest/basic/accessibility.html}{无障碍功能}，这些功能位于 MathJax 上下文菜单中。使用鼠标时，可以在任何公式上点击右键来激活它。屏幕阅读器会提示公式可以被激活（例如，NVDA 在每个公式前会播报“应用程序可点击”以表明存在 MathJax 菜单。）如果您的屏幕阅读器无法读出 $\forall x(A(x) \land B(x))$ 等公式，您可能需要在 MathJax 无障碍菜单的 Speech 子菜单中激活 Speech Output，然后重新加载页面。

以下是一些常见的符号。不同的屏幕阅读器和 TTS 应用程序对它们的发音方式不同，有时甚至完全不发音。每个符号后均附有逻辑学家惯用的读法。
\begin{itemize}
  \item 大写和小写斜体字母，如 $A$（“大写 A”）和 $x$（“小写 x”）
  \item 大写和小写花体字母，如 $\metav{A}$（“花体大写 A”）和 $\metav{x}$（“花体小写 x”）
  \item 逻辑联结词：$\enot$（“非”）、$\eor$（“或”）、$\eand$（“且”）、$\eif$（“仅当”）、$\eiff$（“当且仅当”）、$\ered$（“矛盾”）
  \item 量词：$\forall$（“任意”）、$\exists$（“存在”）
  \item 元理论符号：$\therefore$（“所以”）、$\proves$（“可证”）、$\nproves$（“不可证”）、$\entails$（“蕴涵”）、$\nentails$（“不蕴涵”）
  \item 模态算子：$\ebox$（“必然”）、$\ediamond$（“可能”）
\end{itemize}

屏幕阅读器应当能读出下标，但可能不会提示其为下标。例如，公式 $A_2$ 可能被读作“upper A sub 2”或仅读作“A2”。

表达式 `\blank{}' 和 `\ifeff'（带两个 `f' 的 `if'）在教材中贯穿始终。在 HTML 版本中，`\blank{}' 包含一个不可见的完整拼写单词 `blank'，屏幕阅读器应会将其大声读出。缩写 `\ifeff' 是 `if and only if'（当且仅当）的简写。在 HTML 版本中，它使用不可见的零宽度连字字符编码，这使屏幕阅读器会将 `\ifeff' 读作 `if-eff' 或 `aye-eff-eff'。

\paragraph{证明} \cref{ch.NDTFL,ch.NDFOL,ch.ML} 中的自然演绎证明使用垂直线来指示子证明在何处开始和结束。这些垂直线从子证明的假设行延伸至其最后一行，并显示在任意给定行的行号与公式之间。这使得证明对低视力或全盲用户构成了特殊的挑战。

为了在 HTML 版本中实现这些证明的无障碍访问，证明被编码为表格。表格的每一行有四列：行号、子证明层级指示器、公式以及依据。子证明层级指示器是一个数字，记录当前行包含于多少层嵌套的子证明中。若该行不在任何子证明中，则为 0；若在一层子证明中，则为 1；若在嵌套于另一个子证明内的子证明中，则为 2，依此类推。在播报子证明层级指示器时，屏幕阅读器还应提示前一行是否刚结束了一个子证明，以及新子证明何时开始。表格的表头行和子证明层级指示器被隐藏，以便证明在视觉上与印刷文本保持一致。

以下是一个此类证明的示例：
\begin{fitchproof}
\hypo{wxyz}{(W \eor X) \eor (Y \eor Z)}\PR
\hypo{xy}{X \eif Y}\PR
\hypo{nz}{\enot Z}\PR
\open
	\hypo{wx}{W \eor X}\AS
	\open
		\hypo{w}{W}\AS
		\have{wy1}{W \eor Y}\oi{w}
	\close
	\open
		\hypo{x}{X}\AS
		\have{y1}{Y}\ce{xy, x}
		\have{wy2}{W \eor Y}\oi{y1}
	\close
	\have{wy3}{W \eor Y}\oe{wx, w-wy1, x-wy2}
\close
\open
	\hypo{yz}{Y \eor Z}\AS
	\have{y}{Y}\ds{yz, nz}
	\have{wy}{W \eor Y}\oi{y}
\close
\have{wy4}{W \eor Y}\oe{wxyz, wx-wy3, yz-wy}
\end{fitchproof}
它有 14 行，包含 3 个前提、2 层子证明嵌套和两对相邻的子证明。例如，从第 5 行开始的子证明在第 6 行结束，而第 7 行开始了另一个子证明。因此，第 6 行和第 7 行的子证明层级数相同，但第 6 行和第 7 行位于不同的子证明中。如果您看不到子证明竖线，就必须特别注意子证明层级数的变化，\emph{以及}某个公式何时是假设。屏幕阅读器对第 7 行的播报应大致如下：“7, close subproof, 2, open subproof, upper X, AS.” 在 HTML 版本中，缩写 `\verb|PR|'（代表 `Premise'（前提））和 `\verb|AS|'（代表 `Assumption'（假设））使用 `\verb|abbr|' 标签编码。

请注意，简单的文本转语音应用程序可能完全无法读出包含证明或真值表等复杂表格。

\paragraph{图表}

文中也包含少量图表。在 HTML 版本中，这些图表配有图像描述，例如：
\begin{center}
\begin{tikzpicture}
	\node[circle, grey_shape] (cat1) {A};
	\node[right=10pt of cat1, diamond, phantom_shape] (cat2)  { } ;
	\node[right=10pt of cat2, circle, white_shape] (cat3)  {C} ;
	\node[right=10pt of cat3, white_shape] (cat4)  {D};
\end{tikzpicture}
\bmlDescription{共有三个形状：第一个是一个标有 A 的灰色圆形，一些空白，第二个是一个标有 C 的白色圆形，第四个是一个标有 D 的白色方形。}
\end{center}
此处的图像描述应读作：“共有三个形状：第一个是一个标有 A 的灰色圆形，一些空白，第二个是一个标有 C 的白色圆形，第四个是一个标有 D 的白色方形。”

\paragraph{符号表}

下表列出了文中使用的所有符号、屏幕阅读器最可能的发音方式及其在文中的含义。在最右列，我们提供了使用 ASCII 字符输入这些符号的建议方法，以备无法输入特殊符号的情况（例如在作业或发给教师的电子邮件中）。

\begin{tabular}{lllll}
  \lxBeginTableHead{}符号\lxTableColumnHead{} & 发音\lxTableColumnHead{} & 含义\lxTableColumnHead{} & ASCII 等效输入\lxTableColumnHead{}\\
  \hline\lxEndTableHead
  $\enot$ & 否定号 & 逻辑否定 & \verb+~+ or \verb+-+\\
  $\eor$ & 或 & 逻辑析取 & \verb+\/+\\
  $\eand$ & 且 & 逻辑合取 & \verb+/\+ or \verb+&+\\
  $\eif$ & 右箭头 & 条件 & \verb+->+ or \verb+>+\\
  $\eiff$ & 左右箭头 & 双条件 & \verb+<->+ or \verb+<>+\\
  $\ered$ & 向上钉 & 矛盾 & \verb+_|_+ or \verb+!?+\\
  $\forall$ & 任意 & 全称量词 & \verb|A| or \verb|@|\\
  $\exists$ & 存在 & 存在量词 & \verb|E| or \verb|3|\\
  $\therefore$ & 所以 & 所以 & \verb|:.|\\
  $\proves$ & 右向钉 & 可证 & \verb+|-+\\
  $\nproves$ & 不可证 & 不可证 & \verb+|/-+\\
  $\entails$ & 真 & 蕴涵 & \verb+|=+\\
  $\nentails$ & 不真 & 不蕴涵 & \verb+|/=+\\
  $\ebox$ & 白方块 & 必然 & \verb+[]+\\
  $\ediamond$ & 白菱形 & 可能 & \verb+<>+
\end{tabular}

下标可用下划线表示，例如，$A_2$ 可写为 \verb|A_2|。